\subsection{Encoded decision problem and the meaning of the parameters}\label{sec:decision60}
For the complexity results, an instance is an ordered list of rational catalog levels and nonnegative opening charges, positive rational probabilities summing to one, rational expected caps and realization ceilings, a rational accepted promise and value threshold $z$, an integer command allowance $m\geq1$, and the rational quadratic reward and service coefficients of Section~\ref{sec:model52}. Each rational is represented by a signed binary numerator and a positive binary denominator in lowest terms. Sorting and validation are polynomial-time preprocessing; duplicate catalog entries are consolidated at their least opening charge. The decision language asks whether an original feasible policy using a nonempty book of at most $m$ commands has net value at least $z$. Invalid or infeasible inputs are no-instances. Denote the total encoding length by $L$ and the number of distinct reduced rational caps by $q_b$.

Theorem~\ref{thm:wone59} is a parameterized many-one lower bound for this language with parameter $m+q_b$, already in the common linear-reward, zero-service-cost subclass. Hence it also excludes an algorithm $f(m)L^{O(1)}$ or $f(q_b)L^{O(1)}$, unless FPT equals W[1]: either such algorithm would be fixed-parameter in the sum. It does not assert W[1] membership or completeness, an ETH lower bound on the exact exponent, strong NP-hardness, or a kernel lower bound. The fixed-cap reconstruction and exact-lattice algorithms address different restrictions. In particular, numerical denominators cannot be replaced by their bit lengths in a pseudo-polynomial complexity expression. The positive result in Section~\ref{sec:tariff60} instead removes denominator dependence from the operation count by controlling opening-fee exceptions.
